\documentclass[10pt,a4paper]{amsart}
\usepackage{amsmath,amssymb,amsthm,graphicx}
% Retained preamble for the Stacks Project chapter {CH} (shared project preamble).
% The chapter file pinned-original-source.tex only carries \input{{preamble}} on its first line,
% so no preamble text lives inside the chapter itself.  The Stacks Project uses one single shared
% preamble for every chapter: the snapshot stages it three times (sources/stacks-chow/preamble.tex,
% sources/stacks-etale-cohomology/preamble.tex and focused-seeds/AGP/sources/stacks-snapshot-20261011/
% preamble.tex) and those copies are identical in size.  The block below is copied mechanically,
% line for line, from sources/stacks-chow/preamble.tex lines 14-18 and 154-217: the Xy-pic setup and
% the native theorem declarations together with the mathematical macros.  Nothing was added, removed
% or redefined; the wrapper class and the amsmath/amssymb/amsthm/graphicx base packages are supplied
% by the assembler.
% For commutative diagrams we use Xy-pic
% Chapter: stacks-spaces-more-morphisms.tex -- More on Morphisms of Spaces.
\usepackage[all]{xy}

% We use 2cell for 2-commutative diagrams.
\xyoption{2cell}

% Theorem environments.
%
\theoremstyle{plain}
\newtheorem{theorem}[subsection]{Theorem}
\newtheorem{proposition}[subsection]{Proposition}
\newtheorem{lemma}[subsection]{Lemma}

\theoremstyle{definition}
\newtheorem{definition}[subsection]{Definition}
\newtheorem{example}[subsection]{Example}
\newtheorem{exercise}[subsection]{Exercise}
\newtheorem{situation}[subsection]{Situation}

\theoremstyle{remark}
\newtheorem{remark}[subsection]{Remark}
\newtheorem{remarks}[subsection]{Remarks}

\numberwithin{equation}{subsection}

% Macros
%
\def\lim{\mathop{\mathrm{lim}}\nolimits}
\def\colim{\mathop{\mathrm{colim}}\nolimits}
\def\Spec{\mathop{\mathrm{Spec}}}
\def\Hom{\mathop{\mathrm{Hom}}\nolimits}
\def\Ext{\mathop{\mathrm{Ext}}\nolimits}
\def\SheafHom{\mathop{\mathcal{H}\!\mathit{om}}\nolimits}
\def\SheafExt{\mathop{\mathcal{E}\!\mathit{xt}}\nolimits}
\def\Sch{\mathit{Sch}}
\def\Mor{\mathop{\mathrm{Mor}}\nolimits}
\def\Ob{\mathop{\mathrm{Ob}}\nolimits}
\def\Sh{\mathop{\mathit{Sh}}\nolimits}
\def\NL{\mathop{N\!L}\nolimits}
\def\CH{\mathop{\mathrm{CH}}\nolimits}
\def\proetale{{pro\text{-}\acute{e}tale}}
\def\etale{{\acute{e}tale}}
\def\QCoh{\mathit{QCoh}}
\def\Ker{\mathop{\mathrm{Ker}}}
\def\Im{\mathop{\mathrm{Im}}}
\def\Coker{\mathop{\mathrm{Coker}}}
\def\Coim{\mathop{\mathrm{Coim}}}

% Boxtimes
%
\DeclareMathSymbol{\boxtimes}{\mathbin}{AMSa}{"02}

%
% Macros for moduli stacks/spaces
%
\def\QCohstack{\mathcal{QC}\!\mathit{oh}}
\def\Cohstack{\mathcal{C}\!\mathit{oh}}
\def\Spacesstack{\mathcal{S}\!\mathit{paces}}
\def\Quotfunctor{\mathrm{Quot}}
\def\Hilbfunctor{\mathrm{Hilb}}
\def\Curvesstack{\mathcal{C}\!\mathit{urves}}
\def\Polarizedstack{\mathcal{P}\!\mathit{olarized}}
\def\Complexesstack{\mathcal{C}\!\mathit{omplexes}}
% \Pic is the operator that assigns to X its picard group, usage \Pic(X)
% \Picardstack_{X/B} denotes the Picard stack of X over B
% \Picardfunctor_{X/B} denotes the Picard functor of X over B
\def\Pic{\mathop{\mathrm{Pic}}\nolimits}
\def\Picardstack{\mathcal{P}\!\mathit{ic}}
\def\Picardfunctor{\mathrm{Pic}}
\def\Deformationcategory{\mathcal{D}\!\mathit{ef}}
\usepackage{cleveref}
\crefname{originallemma}{original lemma}{original lemmas}
\Crefname{originallemma}{Original lemma}{Original lemmas}
\makeatletter
\newlabel{spaces-descent-lemma-descending-property-closed-immersion}{{0420}{source}{}{sourcepointer-1}{}}
\newlabel{spaces-descent-lemma-descending-property-closed-immersion@cref}{{[originallemma][0][]0420}{[1][0][]source}{}{}{}}
\makeatother
\title{Being a thickening of algebraic spaces is fpqc local on the target}
\author{The Stacks Project Authors}
\date{stacks-project:spaces-more-morphisms.tex:local-snapshot-20261011; GNU FDL 1.2 or later}
\begin{document}
\maketitle
\noindent\textit{Editorial note.} All delivered mathematical text is the author's own, copied from the pinned Stacks Project chapter file spaces-more-morphisms.tex; no statement or proof step was rewritten, reordered or removed. The author's closing sentence omitting the first order case is reproduced as written.\par\medskip
\noindent\textit{Editorial note.} This card is the generalisation that replaces the scheme-level card STK-morph-descend-property-thick (more-morphisms.tex, tag 0BPF), withdrawn under the situation-A criterion; the withdrawn registration sits in operations/new-material-registrations/withdrawn/ and the job is kept.\par\medskip
\noindent\textit{Editorial note.} Preamble provenance, as required for this source: the first line of the pinned chapter file is \textbackslash{}input\{preamble\}, and the native preamble was not staged with the source. The retained preamble is the project's single Stacks-wide preamble (the same 217-line file used by the already published chow and etale-cohomology cards), copied mechanically line for line from sources/stacks-chow/preamble.tex lines 14-18 (Xy-pic setup) and 154-217 (theorem declarations and mathematical macros). Not one character was changed, no package was added, no macro was defined, and the excerpt serves only as preamble, never as a span coordinate.\par\medskip
\noindent\textit{Allowed standard prerequisites (editorial).}\begin{itemize}
\item Thickening of algebraic spaces, and the reduction from first order to finite order. Same chapter, the definition of thickening and the discussion in the section on thickenings; the statement itself mentions first order thickenings. Supplies the notion whose local nature is asserted, and the reduction principle the author invokes for the first order case.
\item Descent of the property of being a closed immersion for algebraic spaces along an fpqc covering. Descent on Spaces, Lemma labelled spaces-descent-lemma-descending-property-closed-immersion (canonical Stacks tag 0420); invoked at proof lines 1978-1979. Supplies that X -> X' is a closed immersion once its base changes are; the surjectivity step and the assembly into a thickening are retained in the counted span.
\end{itemize}\medskip
\ifdefined\hypertarget\hypertarget{sourcepointer-1}{}\fi
\noindent\textit{Source reference.} Original lemma 0420: sources/stacks-chow/tags.txt, source lines 5311--5311. Cross-chapter standard prerequisite invoked once at proof lines 1978-1979 to descend the closed-immersion property for algebraic spaces; the surjectivity step and the conclusion are retained in the counted span.\par\medskip
\begin{lemma}
\label{lemma-descending-property-thickening}
The property of being a thickening is fpqc local.
Similarly for first order thickenings.
\end{lemma}

\begin{proof}
The statement means the following: Let $S$ be a scheme and let
$X \to X'$ be a morphism of algebraic spaces over $S$.
Let $\{g_i : X'_i \to X'\}$ be an fpqc covering of algebraic spaces
such that the base change $X_i \to X'_i$ is a thickening for all $i$.
Then $X \to X'$ is a thickening. Since the morphisms $g_i$ are jointly
surjective we conclude that $X \to X'$ is surjective. By
Descent on Spaces, Lemma
\ref{spaces-descent-lemma-descending-property-closed-immersion}
we conclude that $X \to X'$ is a closed immersion.
Thus $X \to X'$ is a thickening. We omit the proof in the
case of first order thickenings.
\end{proof}

\end{document}
